\documentclass[12pt,a4paper]{report}

% ─── Packages ──────────────────────────────────────────────────
\usepackage[utf8]{inputenc}
\usepackage[T1]{fontenc}
\usepackage[french]{babel}
\usepackage{geometry}
\usepackage{graphicx}
\usepackage{xcolor}
\usepackage{hyperref}
\usepackage{titlesec}
\usepackage{fancyhdr}
\usepackage{booktabs}
\usepackage{tabularx}
\usepackage{enumitem}
\usepackage{tcolorbox}
\usepackage{mdframed}
\usepackage{fontawesome5}
\usepackage{lmodern}
\usepackage{microtype}
\usepackage{parskip}
\usepackage{float}

% ─── Mise en page ──────────────────────────────────────────────
\geometry{
  top=2.5cm, bottom=2.5cm,
  left=2.8cm, right=2.8cm,
  headheight=15pt
}

% ─── Couleurs ──────────────────────────────────────────────────
\definecolor{primary}{RGB}{30, 58, 95}
\definecolor{accent}{RGB}{37, 99, 235}
\definecolor{success}{RGB}{22, 163, 74}
\definecolor{warning}{RGB}{217, 119, 6}
\definecolor{danger}{RGB}{220, 38, 38}
\definecolor{lightgray}{RGB}{248, 250, 252}
\definecolor{bordergray}{RGB}{226, 232, 240}

% ─── Boîtes d'information ──────────────────────────────────────
\tcbuselibrary{skins, breakable}

\newtcolorbox{infobox}[1][]{
  colback=blue!5!white, colframe=accent,
  fonttitle=\bfseries, title=#1,
  left=6pt, right=6pt, top=4pt, bottom=4pt,
  breakable
}

\newtcolorbox{warningbox}[1][]{
  colback=orange!8!white, colframe=warning,
  fonttitle=\bfseries, title=#1,
  left=6pt, right=6pt, top=4pt, bottom=4pt,
  breakable
}

\newtcolorbox{tipbox}[1][]{
  colback=green!6!white, colframe=success,
  fonttitle=\bfseries, title=#1,
  left=6pt, right=6pt, top=4pt, bottom=4pt,
  breakable
}

% ─── En-tête et pied de page ───────────────────────────────────
\pagestyle{fancy}
\fancyhf{}
\fancyhead[L]{\small\textcolor{primary}{\textbf{École Les Génies}}}
\fancyhead[R]{\small\textcolor{gray}{Cahier d'Utilisation — Portail Numérique}}
\fancyfoot[C]{\small\thepage}
\renewcommand{\headrulewidth}{0.4pt}
\renewcommand{\footrulewidth}{0pt}

% ─── Titres ────────────────────────────────────────────────────
\titleformat{\chapter}[display]
  {\normalfont\Huge\bfseries\color{primary}}
  {\chaptertitlename\ \thechapter}{20pt}{\Huge}
\titleformat{\section}
  {\normalfont\Large\bfseries\color{primary}}
  {\thesection}{1em}{}
\titleformat{\subsection}
  {\normalfont\large\bfseries\color{accent}}
  {\thesubsection}{1em}{}

% ─── Hyperliens ────────────────────────────────────────────────
\hypersetup{
  colorlinks=true,
  linkcolor=accent,
  urlcolor=accent,
  citecolor=accent,
  pdftitle={Cahier d'Utilisation — Portail Numérique École Les Génies},
  pdfauthor={Service Informatique},
}

% ══════════════════════════════════════════════════════════════
\begin{document}

% ─── Page de titre ─────────────────────────────────────────────
\begin{titlepage}
  \centering
  \vspace*{2cm}

  {\color{primary}\rule{\linewidth}{3pt}}
  \vspace{0.5cm}

  {\Huge\bfseries\color{primary} Portail Numérique Scolaire \par}
  \vspace{0.4cm}
  {\LARGE\color{gray} École Les Génies \par}
  \vspace{0.8cm}

  {\color{primary}\rule{\linewidth}{1.5pt}}
  \vspace{1cm}

  {\huge\bfseries Cahier d'Utilisation \par}
  \vspace{0.5cm}
  {\large Manuel de l'utilisateur final \par}

  \vspace{2cm}

  \begin{tabular}{ll}
    \toprule
    \textbf{Document} & Cahier d'utilisation \\
    \textbf{Version}  & 1.0 \\
    \textbf{Date}     & Juillet 2026 \\
    \textbf{Public cible} & Directeurs, Enseignants, Parents \\
    \textbf{Établissement} & École Les Génies — Yaoundé, Cameroun \\
    \bottomrule
  \end{tabular}

  \vfill

  {\color{primary}\rule{\linewidth}{3pt}}
  {\small\color{gray} Document confidentiel — Usage interne uniquement}
\end{titlepage}

% ─── Table des matières ────────────────────────────────────────
\tableofcontents
\clearpage

% ══════════════════════════════════════════════════════════════
\chapter*{Avant-propos}
\addcontentsline{toc}{chapter}{Avant-propos}

Ce document constitue le \textbf{manuel d'utilisation officiel} du Portail Numérique de l'École Les Génies. Il est destiné à l'ensemble des utilisateurs de la plateforme : administrateurs, enseignants et parents d'élèves.

Le portail numérique centralise la gestion scolaire de l'établissement : inscriptions, notes, bulletins, paiements, discipline et communication entre les parties prenantes.

\begin{infobox}[\faInfoCircle\ Comment lire ce document]
Les sections sont organisées par profil utilisateur. Repérez votre rôle dans la table des matières et rendez-vous directement au chapitre correspondant. Les chapitres 1 et 2 sont communs à tous les utilisateurs.
\end{infobox}

% ══════════════════════════════════════════════════════════════
\chapter{Présentation Générale du Portail}

\section{Architecture de l'application}

Le portail numérique est une application web accessible depuis n'importe quel navigateur moderne. Il est composé de trois espaces distincts, chacun adapté à un profil d'utilisateur :

\begin{center}
\begin{tabular}{lll}
\toprule
\textbf{Espace} & \textbf{Port} & \textbf{Public} \\
\midrule
Interface principale (Admin) & 5173 & Administrateurs \\
Portail Enseignant & 3001 & Enseignants \\
Portail Parent     & 3002 & Parents d'élèves \\
\bottomrule
\end{tabular}
\end{center}

\section{Navigateurs supportés}

\begin{itemize}[leftmargin=*]
  \item \textbf{Google Chrome} 110+ (recommandé)
  \item \textbf{Mozilla Firefox} 115+
  \item \textbf{Microsoft Edge} 110+
  \item \textbf{Safari} 16+
\end{itemize}

\begin{warningbox}[\faExclamationTriangle\ Attention]
Internet Explorer n'est pas supporté. L'utilisation sur mobile est possible mais l'expérience optimale est prévue pour écran d'ordinateur.
\end{warningbox}

% ══════════════════════════════════════════════════════════════
\chapter{Accès et Connexion}

\section{Page d'accueil (Landing Page)}

Avant de se connecter, les utilisateurs arrivent sur la \textbf{page d'accueil publique} de l'école. Elle présente :
\begin{itemize}[leftmargin=*]
  \item Les informations générales de l'établissement
  \item La galerie photo de la vie scolaire
  \item Les témoignages de parents
  \item Un bouton \textbf{«~Portail Numérique~»} pour accéder à la connexion
\end{itemize}

\section{Écran de connexion}

L'écran de connexion est le point d'entrée unique pour tous les utilisateurs de la plateforme :

\begin{enumerate}[leftmargin=*]
  \item Saisir son \textbf{identifiant} et son \textbf{mot de passe}
  \item Cliquer sur \textbf{«~Se connecter~»}
  \item Le système détecte automatiquement votre profil (Administrateur, Enseignant, ou Parent) et vous redirige vers le tableau de bord correspondant.
\end{enumerate}

\begin{warningbox}[\faLock\ Sécurité]
Ne communiquez jamais votre mot de passe à des tiers. En cas de perte, contactez le service informatique pour obtenir une réinitialisation.
\end{warningbox}

\section{Déconnexion}

La déconnexion est accessible depuis le menu latéral gauche, en bas, via l'icône de sortie. Il est recommandé de toujours se déconnecter après utilisation sur un poste partagé.

% ══════════════════════════════════════════════════════════════
\chapter{Espace Administrateur}

L'espace administrateur est le cœur de gestion de l'établissement. Il est accessible aux directeurs, fondateurs, intendants et membres de l'administration.

\section{Tableau de bord}

À la connexion, l'administrateur accède au \textbf{tableau de bord} qui affiche une vue d'ensemble :
\begin{itemize}[leftmargin=*]
  \item Nombre total d'élèves inscrits
  \item Nombre d'enseignants
  \item Paiements en attente (montant global)
  \item Taux de présence du jour
  \item Graphiques d'évolution des effectifs et des paiements
  \item Dernières activités du système
\end{itemize}

\begin{figure}[H]
    \centering
    \includegraphics[width=0.9\textwidth]{"images/Screenshot 2026-07-09 at 09-51-25 Application de gestion scolaire.png"}
    \caption{Tableau de bord de l'administrateur}
\end{figure}

\section{Gestion Académique}

\subsection{Classes et Cycles}
\textit{Menu : Gestion académique > Cycles \& Classes}

Permet de créer et gérer les niveaux scolaires (Maternelle, CP, CE1, CE2, CM1, CM2) ainsi que les classes correspondantes.

\textbf{Opérations disponibles :}
\begin{itemize}[leftmargin=*]
  \item Ajouter un nouveau cycle ou une nouvelle classe
  \item Modifier le libellé d'une classe existante
  \item Supprimer une classe (si aucun élève n'y est inscrit)
\end{itemize}

\begin{figure}[H]
    \centering
    \includegraphics[width=0.9\textwidth]{"images/Screenshot 2026-07-09 at 09-54-52 Application de gestion scolaire.png"}
    \caption{Gestion des classes et des cycles}
\end{figure}

\subsection{Salles}
\textit{Menu : Gestion académique > Salles}

Permet de gérer les salles de classe physiques de l'établissement. Chaque salle est associée à une classe et peut recevoir un titulaire.

\begin{figure}[H]
    \centering
    \includegraphics[width=0.9\textwidth]{"images/Screenshot 2026-07-09 at 09-55-01 Application de gestion scolaire.png"}
    \caption{Gestion des salles}
\end{figure}

\subsection{Titulaires}
\textit{Menu : Gestion académique > Titulaires}

Permet d'affecter un enseignant titulaire à une salle/classe. Un titulaire est l'enseignant principal responsable d'un groupe d'élèves.

\textbf{Pour affecter un titulaire :}
\begin{enumerate}[leftmargin=*]
  \item Cliquer sur «~Ajouter un titulaire~»
  \item Sélectionner l'enseignant dans la liste déroulante
  \item Sélectionner la salle concernée
  \item Valider
\end{enumerate}
\begin{figure}[H]
    \centering
    \includegraphics[width=0.9\textwidth]{"images/Screenshot 2026-07-09 at 09-55-12 Application de gestion scolaire.png"}
    \caption{Affectation des titulaires}
\end{figure}

\section{Gestion des Élèves}

\subsection{Liste des élèves}
\textit{Menu : Élèves > Liste des élèves}

Affiche l'ensemble des élèves inscrits avec leurs informations : nom, prénom, matricule, classe, statut, moyenne générale, absences.

\textbf{Fonctionnalités de recherche :}
\begin{itemize}[leftmargin=*]
  \item Barre de recherche par nom, prénom ou matricule
  \item Filtrage par classe ou statut
  \item Clic sur un élève pour voir son profil complet
\end{itemize}

\begin{figure}[H]
    \centering
    \includegraphics[width=0.9\textwidth]{"images/Screenshot 2026-07-09 at 09-52-55 Application de gestion scolaire.png"}
    \caption{Liste globale des élèves}
\end{figure}

\begin{figure}[H]
    \centering
    \includegraphics[width=0.9\textwidth]{"images/Screenshot 2026-07-09 at 09-53-32 Application de gestion scolaire.png"}
    \caption{Profil détaillé d'un élève}
\end{figure}

\subsection{Inscriptions}
\textit{Menu : Élèves > Inscriptions}

Permet d'inscrire un nouvel élève dans l'établissement.

\textbf{Processus d'inscription :}
\begin{enumerate}[leftmargin=*]
  \item Cliquer sur «~Nouvelle inscription~»
  \item Remplir le formulaire : état civil de l'élève (nom, prénom, date et lieu de naissance, sexe)
  \item Renseigner les informations des parents (nom, contact)
  \item Sélectionner la classe d'affectation
  \item Valider l'inscription — un matricule unique est généré automatiquement
\end{enumerate}

\begin{figure}[H]
    \centering
    \includegraphics[width=0.9\textwidth]{"images/Screenshot 2026-07-09 at 09-51-58 Application de gestion scolaire.png"}
    \caption{Formulaire de pré-inscription (Étape 1)}
\end{figure}

\begin{figure}[H]
    \centering
    \includegraphics[width=0.9\textwidth]{"images/Screenshot 2026-07-09 at 09-51-58 Application de gestion scolaire.png"}
    \caption{Validation de l'inscription}
\end{figure}

\begin{tipbox}[\faLightbulb\ Conseil]
Vérifiez qu'un élève n'est pas déjà inscrit en le recherchant par son nom avant de créer une nouvelle fiche, afin d'éviter les doublons.
\end{tipbox}

\section{Pédagogie}

\subsection{Cours et Matières}
\textit{Menu : Pédagogie > Cours \& Matières}

Permet de créer et gérer le catalogue des matières enseignées, avec leur coefficient respectif.

\begin{figure}[H]
    \centering
    \includegraphics[width=0.9\textwidth]{"images/Screenshot 2026-07-09 at 10-20-54 Application de gestion scolaire.png"}
    \caption{Catalogue des matières}
\end{figure}

\subsection{Emploi du Temps}
\textit{Menu : Pédagogie > Emploi du temps}

Permet de visualiser et modifier l'emploi du temps par classe et par jour de la semaine.

\begin{figure}[H]
    \centering
    \includegraphics[width=0.9\textwidth]{"images/Screenshot 2026-07-09 at 11-19-07 Application de gestion scolaire.png"}
    \caption{Emploi du temps}
\end{figure}


\subsection{Épreuves}
\textit{Menu : Pédagogie > Épreuves}

Permet de planifier les évaluations (contrôles, examens) en indiquant la matière, la date et le type d'épreuve.

\begin{figure}[H]
    \centering
    \includegraphics[width=0.9\textwidth]{"images/Screenshot 2026-07-09 at 11-21-27 Application de gestion scolaire.png"}
    \caption{Epreuves}
\end{figure}


\subsection{Notes et Moyennes}
\textit{Menu : Pédagogie > Notes \& Moyennes}

Affiche un récapitulatif des performances académiques de tous les élèves :
\begin{itemize}[leftmargin=*]
  \item Moyenne générale de l'école
  \item Taux de réussite global
  \item Liste des élèves avec leur moyenne individuelle
  \item Identification des élèves en difficulté (moyenne $<$ 10/20)
\end{itemize}

\begin{figure}[H]
    \centering
    \includegraphics[width=0.9\textwidth]{"images/Screenshot 2026-07-09 at 11-23-36 Application de gestion scolaire.png"}
    \caption{Notes \& Moyennes}
\end{figure}

\subsection{Bulletins PDF}
\textit{Menu : Pédagogie > Bulletins PDF}

Permet de générer et télécharger les bulletins scolaires au format PDF pour n'importe quel trimestre.

\begin{figure}[H]
    \centering
    \includegraphics[width=0.9\textwidth]{"images/Screenshot 2026-07-09 at 09-55-37 Application de gestion scolaire.png"}
    \caption{Bulletins PDF}
\end{figure}

\section{Suivi des Élèves}

\subsection{Présences}
\textit{Menu : Suivi > Présences}

Permet de consulter le registre des absences et des retards, filtré par date ou par élève.

\begin{figure}[H]
    \centering
    \includegraphics[width=0.9\textwidth]{"images/Screenshot 2026-07-09 at 11-27-36 Application de gestion scolaire.png"}
    \caption{Presences}
\end{figure}


\subsection{Discipline}
\textit{Menu : Suivi > Discipline}

Permet de saisir et consulter les sanctions disciplinaires. Chaque sanction comporte :
\begin{itemize}[leftmargin=*]
  \item L'élève concerné
  \item La date de l'incident
  \item Le motif de la sanction
  \item Le nombre de points de conduite retirés
\end{itemize}

\begin{figure}[H]
    \centering
    \includegraphics[width=0.9\textwidth]{"images/Screenshot 2026-07-09 at 11-27-44 Application de gestion scolaire.png"}
    \caption{Discipline}
\end{figure}


\section{Gestion Financière}

\subsection{Paiements}
\textit{Menu : Financier > Paiements}

Affiche l'historique de tous les paiements reçus (frais de scolarité, cantines, etc.) avec les détails de chaque transaction.

\begin{figure}[H]
    \centering
    \includegraphics[width=0.9\textwidth]{"images/Screenshot 2026-07-09 at 09-54-19 Application de gestion scolaire.png"}
    \caption{Historique des paiements}
\end{figure}

\subsection{Impayés}
\textit{Menu : Financier > Impayés}

Liste les familles ayant des paiements en retard, avec le montant dû et la date d'échéance dépassée.

\begin{figure}[H]
    \centering
    \includegraphics[width=0.9\textwidth]{"images/Screenshot 2026-07-09 at 11-34-21 Application de gestion scolaire.png"}
    \caption{Impayes}
\end{figure}
\section{Communication}

\subsection{Messagerie}
\textit{Menu : Communication > Messagerie}

Permet d'envoyer des messages aux parents d'élèves individuellement ou en groupe.

\textbf{Pour envoyer un message :}
\begin{enumerate}[leftmargin=*]
  \item Cliquer sur «~Nouveau message~»
  \item Sélectionner le ou les destinataires (parent ou classe entière)
  \item Saisir l'objet et le contenu du message
  \item Envoyer
\end{enumerate}

\section{Administration Système}

\subsection{Gestion des Utilisateurs}
\textit{Menu : Administration > Utilisateurs}

Permet de créer des comptes pour les enseignants et les parents, de réinitialiser les mots de passe et de gérer les droits d'accès.

\textbf{Créer un compte utilisateur :}
\begin{enumerate}[leftmargin=*]
  \item Cliquer sur «~Nouvel utilisateur~»
  \item Choisir le rôle (Enseignant ou Parent)
  \item Renseigner les informations personnelles et l'identifiant de connexion
  \item Un mot de passe temporaire est généré — à communiquer à l'utilisateur
\end{enumerate}

\begin{figure}[H]
    \centering
    \includegraphics[width=0.9\textwidth]{"images/Screenshot 2026-07-09 at 09-54-19 Application de gestion scolaire.png"}
    \caption{Liste des utilisateurs du système}
\end{figure}

\begin{figure}[H]
    \centering
    \includegraphics[width=0.9\textwidth]{"images/Screenshot 2026-07-09 at 09-54-31 Application de gestion scolaire.png"}
    \caption{Création d'un nouveau compte}
\end{figure}

\subsection{Configuration}
\textit{Menu : Administration > Configuration}


% ══════════════════════════════════════════════════════════════
\chapter{Espace Enseignant}

L'espace enseignant est accessible sur le port 3001. Il est optimisé pour une utilisation quotidienne en classe.

\section{Tableau de bord Enseignant}

À la connexion, le tableau de bord de l'enseignant affiche :
\begin{itemize}[leftmargin=*]
  \item Sa classe assignée et le nombre d'élèves
  \item La répartition filles / garçons
  \item Les statistiques d'absences
  \item La moyenne générale de la classe
  \item Les dernières notes saisies
\end{itemize}

\begin{figure}[H]
    \centering
    \includegraphics[width=0.9\textwidth]{"images/Screenshot 2026-07-09 at 09-56-56 Espace Enseignant Les Genies.png"}
    \caption{Tableau de bord de l'enseignant}
\end{figure}

\begin{infobox}[\faInfoCircle\ Information]
Si le message «~Aucune classe assignée~» s'affiche, contactez l'administrateur pour qu'il vous affecte à une classe via le menu Titulaires.
\end{infobox}

\section{Gestion des Élèves (vue Enseignant)}
\textit{Menu : Mes élèves}

Affiche la liste complète des élèves de la classe avec :
\begin{itemize}[leftmargin=*]
  \item La photo de profil (avatar coloré selon le genre)
  \item La moyenne individuelle avec indicateur de tendance
  \item Le nombre d'absences et de retards
  \item Un bouton \textbf{«~Voir profil~»} pour afficher la fiche détaillée
\end{itemize}

\subsection{Profil détaillé d'un élève}
En cliquant sur «~Voir profil~», un panneau latéral s'ouvre avec :
\begin{itemize}[leftmargin=*]
  \item L'état civil complet (date et lieu de naissance)
  \item La moyenne et le statut académique
  \item Le bilan absences/retards
\end{itemize}

\section{Gestion des Absences}
\textit{Menu : Absences}

Permet de saisir les absences et retards du jour pour les élèves de la classe.

\textbf{Saisir une absence :}
\begin{enumerate}[leftmargin=*]
  \item Se rendre sur la page «~Absences~»
  \item Sélectionner la date (par défaut : aujourd'hui)
  \item Cocher les élèves absents dans la liste
  \item Indiquer si c'est une absence ou un retard
  \item Enregistrer
\end{enumerate}

\begin{figure}[H]
    \centering
    \includegraphics[width=0.9\textwidth]{"images/Screenshot 2026-07-09 at 09-57-11 Espace Enseignant Les Genies.png"}
    \caption{Saisie de l'appel et des présences}
\end{figure}

\section{Évaluations et Notes}

\subsection{Créer une Épreuve (Vue à développer)}
\textit{Emplacement prévu : Menu > Évaluations > Nouvelle Épreuve}

Cette nouvelle interface permettra à l'enseignant de créer une épreuve, d'en définir la matière, le trimestre, le type (Séquence, Examen) et le coefficient. Une fois l'épreuve créée, elle s'ajoutera à la liste pour permettre la saisie des notes.

\begin{figure}[H]
    \centering
    \includegraphics[width=0.9\textwidth]{"images/Screenshot 2026-07-09 at 09-57-52 Espace Enseignant Les Genies.png"}
    \caption{Creer une epreuve}
\end{figure}

\subsection{Saisie des Notes}
\textit{Menu : Notes}

Permet de saisir les notes des élèves pour une épreuve donnée, avec un tableau listant tous les élèves de la classe.

\begin{figure}[H]
    \centering
    \includegraphics[width=0.9\textwidth]{"images/Screenshot 2026-07-09 at 09-57-26 Espace Enseignant Les Genies.png"}
    \caption{Saisie des notes}
\end{figure}

\section{Organisation Pédagogique}


\subsection{Emploi du Temps}
\textit{Menu : Emploi du temps}

Affiche la grille horaire hebdomadaire de l'enseignant.

\begin{figure}[H]
    \centering
    \includegraphics[width=0.9\textwidth]{"images/Screenshot 2026-07-09 at 12-00-29 Espace Enseignant Les Genies.png"}
    \caption{Emploi du temps}
\end{figure}

\section{Consulter les Bulletins}
\textit{Menu : Bulletins PDF}

L'enseignant peut avoir un aperçu des bulletins trimestriels générés pour sa classe et les exporter au format PDF.

% ══════════════════════════════════════════════════════════════
\chapter{Espace Parent}

L'espace parent est accessible sur le port 3002. Il permet aux parents de suivre la scolarité de leurs enfants à distance.

\section{Tableau de bord Parent}

Après connexion, le parent accède à un tableau de bord personnalisé présentant :

\begin{itemize}[leftmargin=*]
  \item La \textbf{liste de ses enfants} inscrits, avec pour chacun :
    \begin{itemize}
      \item Le nom de la classe et le nom du professeur titulaire
      \item La moyenne générale et le statut (Excellent / Bien / À surveiller)
    \end{itemize}
  \item Les \textbf{statistiques rapides} (nombre d'enfants, bulletins disponibles, paiements)
  \item Le \textbf{suivi financier} (paiements à effectuer)
  \item La \textbf{messagerie} (messages reçus de l'administration)
\end{itemize}

\begin{figure}[H]
    \centering
    \includegraphics[width=0.9\textwidth]{"images/Screenshot 2026-07-09 at 09-58-34 Les Genies - Espace Parent.png"}
    \caption{Tableau de bord du parent}
\end{figure}

\section{Consulter les Notes}
\textit{Menu : Notes}

La page «~Carnet de Notes~» affiche, pour chaque enfant, l'ensemble des notes obtenues dans chaque matière :

\begin{itemize}[leftmargin=*]
  \item Organisation par matière avec le coefficient
  \item Chaque évaluation avec sa note et sa date
  \item La moyenne par matière avec code couleur :
    \begin{itemize}
      \item \textcolor{success}{\textbf{Vert}} : Note $\geq$ 16/20
      \item \textcolor{accent}{\textbf{Bleu}} : Note $\geq$ 12/20
      \item \textcolor{warning}{\textbf{Orange}} : Note $\geq$ 10/20
      \item \textcolor{danger}{\textbf{Rouge}} : Note $<$ 10/20
    \end{itemize}
\end{itemize}

Si le parent a plusieurs enfants, des boutons permettent de basculer d'un enfant à l'autre.

\begin{figure}[H]
    \centering
    \includegraphics[width=0.9\textwidth]{"images/Screenshot 2026-07-09 at 09-59-29 Les Genies - Espace Parent.png"}
    \caption{Suivi académique (Notes)}
\end{figure}

\section{Consulter les Bulletins}
\textit{Menu : Bulletins}

\subsection{Liste des bulletins disponibles}
La page affiche tous les bulletins disponibles par trimestre et par enfant. Chaque carte de bulletin indique :
\begin{itemize}[leftmargin=*]
  \item Le nom de l'enfant et le trimestre
  \item La moyenne générale du trimestre
  \item La mention obtenue
  \item Le rang dans la classe
\end{itemize}

\begin{figure}[H]
    \centering
    \includegraphics[width=0.9\textwidth]{"images/Screenshot 2026-07-09 at 09-58-51 Les Genies - Espace Parent.png"}
    \caption{Liste des bulletins disponibles}
\end{figure}

\subsection{Ouvrir un bulletin}
\begin{enumerate}[leftmargin=*]
  \item Cliquer sur la carte d'un bulletin
  \item Le bulletin officiel s'ouvre en plein écran au format A4
\end{enumerate}

\subsection{Contenu du bulletin officiel}
Le bulletin affiché contient toutes les informations officielles :
\begin{itemize}[leftmargin=*]
  \item \textbf{En-tête} : nom de l'école, République du Cameroun, trimestre et année académique
  \item \textbf{Identité de l'élève} : nom, prénom, matricule, date et lieu de naissance, classe
  \item \textbf{Tableau des notes} : matière, coefficient, notes obtenues, moyenne et appréciation par matière
  \item \textbf{Résumé} : moyenne générale, rang, mention, absences, retards
  \item \textbf{Conduite et discipline} : récapitulatif des éventuelles sanctions du trimestre
  \item \textbf{Appréciation générale} du maître/de la maîtresse
  \item \textbf{Zones de signature} : Directeur, Maître/Maîtresse, Parent
\end{itemize}

\begin{figure}[H]
    \centering
    \includegraphics[width=0.9\textwidth]{"images/Screenshot 2026-07-09 at 09-59-11 Les Genies - Espace Parent.png"}
    \caption{Vue Bulletin}
\end{figure}


\subsection{Imprimer ou sauvegarder en PDF}
\begin{enumerate}[leftmargin=*]
  \item Cliquer sur le bouton bleu \textbf{«~Imprimer / PDF~»} en haut à droite
  \item Une fenêtre d'impression du navigateur s'ouvre
  \item Pour sauvegarder en PDF : choisir «~Enregistrer en PDF~» comme imprimante
  \item Le document est formaté pour une impression A4 en portrait
\end{enumerate}

\begin{tipbox}[\faLightbulb\ Conseil impression]
Dans les paramètres d'impression, activez «~Graphiques d'arrière-plan~» pour que les couleurs et les bordures du bulletin soient correctement imprimées.
\end{tipbox}

\section{Messagerie Parent}
\textit{Menu : Messages}

Permet de consulter les messages envoyés par l'administration ou les enseignants. Les messages non lus sont mis en évidence.

\section{Paiements}
\textit{Menu : Paiements}

Affiche l'historique des paiements effectués pour les frais de scolarité et les éventuels montants restants à régler.

\begin{figure}[H]
    \centering
    \includegraphics[width=0.9\textwidth]{"images/Screenshot 2026-07-09 at 09-59-41 Les Genies - Espace Parent.png"}
    \caption{Suivi financier du parent}
\end{figure}

% ══════════════════════════════════════════════════════════════
\chapter{Questions Fréquentes}

\section{Problèmes de connexion}

\begin{tabularx}{\textwidth}{lX}
\toprule
\textbf{Problème} & \textbf{Solution} \\
\midrule
Identifiant ou mot de passe incorrect & Vérifiez les majuscules/minuscules. Contactez l'administration pour réinitialisation. \\
\addlinespace
Page blanche ou erreur au chargement & Vérifiez votre connexion internet. Attendez quelques secondes et actualisez la page (F5). \\
\addlinespace
Le bulletin ne s'affiche pas & Vérifiez que des notes ont bien été saisies pour ce trimestre. Contactez l'administration. \\
\addlinespace
Mon enfant n'apparaît pas & Vérifiez que l'inscription a bien été validée par l'administration. \\
\bottomrule
\end{tabularx}

\section{Problèmes d'impression}

\begin{tabularx}{\textwidth}{lX}
\toprule
\textbf{Problème} & \textbf{Solution} \\
\midrule
Le bulletin s'imprime sans couleurs & Dans les options d'impression, activer «~Graphiques d'arrière-plan~». \\
\addlinespace
Le tableau déborde sur la page & Vérifier que le format de page est bien A4 et l'orientation Portrait. \\
\addlinespace
La fenêtre d'impression ne s'ouvre pas & Désactiver le bloqueur de popups du navigateur pour ce site. \\
\bottomrule
\end{tabularx}

% ══════════════════════════════════════════════════════════════
\chapter{Contacts et Support}

En cas de difficultés techniques, veuillez contacter le service informatique de l'établissement :

\begin{center}
\begin{tabular}{ll}
\toprule
\textbf{Service} & \textbf{Contact} \\
\midrule
Support technique & informatique@lesgenies.cm \\
Direction        & direction@lesgenies.cm \\
Téléphone        & (+237) xxx xxx xxx \\
\bottomrule
\end{tabular}
\end{center}

\vspace{1cm}

\begin{infobox}[\faInfoCircle\ Signaler un bug]
Pour signaler un dysfonctionnement, notez précisément : le message d'erreur affiché, les actions effectuées avant l'erreur, votre navigateur et l'heure de l'incident. Ces informations permettront une résolution rapide.
\end{infobox}

% ─── Fin de document ───────────────────────────────────────────
\vfill
\begin{center}
  \color{gray}\small
  \textit{École Les Génies — Portail Numérique v1.0 — Juillet 2026} \\
  \textit{Document interne — Reproduction interdite sans autorisation}
\end{center}

\end{document}
